\section{Exponential errors and inverse polynomial mass}
\label{sec:exponential_polynomial_absorption}

The following scalar estimate is used to compare a finite sum of
exponential errors with an inverse polynomial mass in
\cite[Section~7]{OpenAI2026AreaLaw}.

\begin{theorem}[Exponential decay below a shifted inverse power]
    \label{thm:exponential_polynomial_absorption}
    \lean{Real.eventually_mul_exp_neg_le_inv_four_pow_add_one}
    \leanok
    Let $c>0$, let $m\in\mathbb R$, and let $d$ be a nonnegative integer.
    For all sufficiently large integers $k$,
    \begin{align*}
        m e^{-ck} & \le\frac{1}{4(k+1)^d}.
    \end{align*}
    In particular, the statement applies when $m$ counts a finite family of
    cutoffs, including the empty family.
\end{theorem}
\begin{proof}\leanok
    For fixed $c>0$ and $d\ge0$, the quantity $x^d e^{-cx}$ tends to zero
    as $x\to\infty$. Consequently
    \begin{align*}
        4m(k+1)^d e^{-ck}
         & =4m e^c\,(k+1)^d e^{-c(k+1)}
        \longrightarrow 0.
    \end{align*}
    It is therefore eventually at most one. Division by the positive
    quantity $4(k+1)^d$ gives the conclusion.
\end{proof}

\section{Finitely many exponential losses}
\label{sec:finite_exponential_retained_mass}

The following scalar estimates compare finitely many regional failures
with an inverse polynomial label mass, as required in
\cite[Section~7]{OpenAI2026AreaLaw}. They do not assume that the
individual exponential rates are equal.

\begin{theorem}[A finite sum of exponential errors]
\label{thm:finite_exponential_polynomial_absorption}
\lean{Real.eventually_sum_exp_neg_le_inv_four_pow_add_one}
    \leanok
\uses{thm:exponential_polynomial_absorption}
Let $I$ be a finite set, let $c_i>0$ for every $i\in I$, and let
$d\ge0$ be an integer. For all sufficiently large integers $k$,
\begin{align}
 \sum_{i\in I}e^{-c_i k}&\le\frac{1}{4(k+1)^d}.
\end{align}
The set $I$ may be empty.
\end{theorem}
\begin{proof}
    \leanok
If $I$ is empty, the left side is zero. Otherwise write $m=|I|>0$.
For each $i$, exponential decay gives
$m e^{-c_i k}\le 1/(4(k+1)^d)$ eventually. Finiteness gives a common
threshold after which all these inequalities hold. Sum them and cancel
the positive factor $m$.
\end{proof}

\begin{theorem}[Retention of half the original mass]
\label{thm:half_mass_after_finite_exponential_losses}
\lean{Real.eventually_half_mass_of_exponential_losses}
    \leanok
\uses{thm:finite_exponential_polynomial_absorption}
Let $q_k,r_k$ and $\ell_{i,k}$ be real sequences indexed by integers
$k\ge0$, with $i$ in a finite set $I$, and let $d\ge0$ be an integer.
Suppose $q_k\ge 1/(2(k+1)^d)$ for all sufficiently large $k$ and
$r_k\ge q_k-\sum_{i\in I}\ell_{i,k}$ for every $k$.
Suppose also that, for each $i\in I$, there is a constant $c_i>0$ for which
$\ell_{i,k}\le e^{-c_i k}$ for all sufficiently large $k$.
Then, for all sufficiently large $k$,
\begin{align}
 r_k&\ge\frac{q_k}{2},
 &r_k&\ge\frac{1}{4(k+1)^d}.
\end{align}
No additional sign assumptions on the sequences are required.
\end{theorem}
\begin{proof}
    \leanok
Intersect the finitely many eventual tail bounds and apply Theorem~\ref{thm:finite_exponential_polynomial_absorption}. Thus, with $t_k=1/(4(k+1)^d)$, the total loss is eventually at
most $t_k$, while $q_k\ge2t_k$. The retained-mass inequality gives
\begin{align}
 r_k&\ge q_k-t_k\ge q_k/2\ge t_k.
\end{align}
The quantity $q_k$ is the same original mass; no new selection is involved.
\end{proof}

\section{Functional calculus of the finite-spectrum surprisal}
\label{sec:surprisal_functional_calculus}

For a real number $u$, put $\ell(u)=\log|u|$ when $u\ne0$ and
$\ell(0)=0$. For a finite Hermitian matrix $A$, define $\ell(A)$ by
spectral functional calculus. On positive eigenvalues this is the usual
logarithm, with value zero on the kernel. This convention allows the
two cutoff expressions in \cite[Section~7]{OpenAI2026AreaLaw} to be
identified also for singular densities.

\begin{theorem}[Composition with the negative logarithm]
    \label{thm:surprisal_functional_calculus}
    \lean{Matrix.IsHermitian.cfc_neg_log_eq}
    \leanok
    For every finite Hermitian matrix $A$ and every function
    $f:\mathbb R\to\mathbb R$,
    \begin{align*}
        f(-\ell(A)) & =(f\circ(-\ell))(A).
    \end{align*}
    No positivity or invertibility of $A$, and no continuity of $f$ on the
    whole real line, is required. In particular, for any real $a$, the strict
    upper spectral cutoff of $-\ell(A)$ equals the functional calculus on
    $A$ of the function $u\mapsto\mathbf1_{\{a<-\ell(u)\}}$.
\end{theorem}
\begin{proof}\leanok
    Every function on the finite spectrum of $A$ is continuous in its
    relative topology, and the same holds on its image under $-\ell$.
    The composition rule for spectral functional calculus therefore applies.
    At a zero eigenvalue, both sides have the value $f(0)$.
\end{proof}

\section{Exponential upper tails of finite-copy surprisal}
\label{sec:iid_surprisal_exponential}

The concentration step in \cite[Section~7]{OpenAI2026AreaLaw} concerns a
fixed finite probability law. Its exponential rate may depend on that law.
Throughout, $0\log 0=0$; the value assigned to the logarithm at a zero
weight does not affect any probability or trace below.

\begin{theorem}[The upper tail of a finite product law]
    \label{thm:iid_surprisal_exponential_finite_law}
    \lean{Entropy.exists_surprisal_upper_tail_pi_le_exp}
    \leanok
    Let $(p_i)_{i\in I}$ be a finite family with $p_i\ge0$ and
    $\sum_i p_i=1$, and set $H=-\sum_i p_i\log p_i$. There is a constant
    $c>0$, depending only on this family, such that for every integer $k\ge1$,
    \begin{align*}
        \sum_{x\in I^k:\,-\log\prod_{j=0}^{k-1} p_{x_j}>k(H+1)}
        \prod_{j=0}^{k-1} p_{x_j}
         & \le e^{-ck}.
    \end{align*}
    Tuples containing a zero weight contribute zero to the sum.
\end{theorem}
\begin{proof}\leanok
    Define $X(i)=-\log p_i$ when $p_i>0$, and $X(i)=0$ otherwise. Since
    $p_i\le1$, the variable $X$ lies in $[0,M]$, where
    $M=1+\sum_i|X(i)|>0$. Its expectation under the law $(p_i)_i$ is $H$.
    For independent samples $X_0,\ldots,X_{k-1}$, Hoeffding's inequality gives
    \begin{align*}
        \mathbb P\!\left(\sum_{j=0}^{k-1}(X_j-H)\ge k\right)
         & \le \exp\!\left(-\frac{2k}{M^2}\right).
    \end{align*}
    On every tuple of positive probability, the sum of the surprisals equals
    the surprisal of the product weight. The strict event in the statement is
    contained in the displayed event. Thus $c=2/M^2$ is admissible, including
    when the probability law is concentrated at a single point.
\end{proof}

\begin{theorem}[The upper tail of a density matrix]
    \label{thm:iid_surprisal_exponential_density}
    \lean{Matrix.PosSemidef.exists_re_trace_surprisal_upper_tail_finKronecker_le_exp}
    \leanok
    \uses{thm:iid_surprisal_exponential_finite_law,thm:iid_surprisal_spectral_distribution}
    Let $\rho\ge0$ be a matrix on a finite-dimensional Hilbert space with
    $\operatorname{Tr}\rho=1$. Write $S=S(\rho)$,
    $\rho_k=\rho^{\otimes k}$, and $L_k=-\log\rho_k$, taking the logarithm
    to be zero on the kernel. There is a constant $c>0$, depending only on
    $\rho$, such that for every integer $k\ge1$,
    \begin{align*}
        \operatorname{Re}\operatorname{Tr}\!\left(
        \rho_k\mathbf1_{(k(S+1),\infty)}(L_k)\right)
         & \le e^{-ck}.
    \end{align*}
    The matrix $\rho$ may be singular.
\end{theorem}
\begin{proof}\leanok
    The eigenvalues of $\rho$ form a finite probability distribution whose
    entropy is $S$. The product spectral distribution identifies the trace
    with the preceding finite-product tail. Zero eigenvalues have zero trace
    weight, so the value of the logarithm on the kernel causes no additional
    term.
\end{proof}

\section{A common upper cutoff for the Schur labels}
\label{sec:schur-common-cutoffs}

Let $F$ be a finite set of sites with finite coordinate sets $I_f$,
and let $k\ge0$. For a region $B\subseteq F$, write
$P_B^\lambda$ for the central isotypic projection of its copy-permutation
action and
\begin{align}
 C_B(a)&=\sum_{\lambda:\,\log\dim\lambda\le a}P_B^\lambda.
\end{align}
Thus $C_B(a)$ is the closed cutoff of
Definition~\ref{def:schur_label_cutoff}. The following simultaneous
projection estimates are used in
\cite[Section~7, lines~332--354]{OpenAI2026AreaLaw}.

\begin{definition}[Common cutoff]
\label{def:schur-common-cutoff}
\lean{TensorPower.commonLabelCutoff}
    \leanok
\uses{def:schur_label_cutoff}
For a finite list $\mathcal B=(B_1,\ldots,B_m)$ and real thresholds
$a_B$, put $C_j=C_{B_j}(a_{B_j})$ and set
$C=\prod_{j=1}^m C_j$, in the order of the list.
The empty product is the identity. Repeated regions are permitted.
\end{definition}

\begin{theorem}[Common projection and finite union bound]
\label{thm:schur-common-cutoff-mass}
\lean{TensorPower.norm_commonLabelCutoff_auxLabel_ge}
    \leanok
\uses{def:schur-common-cutoff,thm:schur_label_commutation}
Suppose that the regions in $\mathcal B$ are pairwise comparable by
inclusion, and let $R\subseteq F$ be disjoint from every $B_j$.
For any irreducible label $\lambda$ of $\mathfrak S_k$, put
$Q=P_R^\lambda$. Then every vector $\psi$ satisfies
\begin{align}
 \|CQ\psi\|^2
 &\ge \|Q\psi\|^2-\sum_{j=1}^m\|(I-C_j)\psi\|^2.
\end{align}
\end{theorem}
\begin{proof}
    \leanok
Nestedness makes the $C_j$ commute, and disjointness makes each $C_j$
commute with $Q$. For an orthogonal projection $P$,
$\|Pv\|^2+\|(I-P)v\|^2=\|v\|^2$.
Apply this identity successively. At each step the removed component is
a commuting orthogonal projection of $(I-C_j)\psi$, so its norm is
at most $\|(I-C_j)\psi\|$. Summing the losses gives the inequality.
\end{proof}

\begin{theorem}[Normalized symmetric common projection]
\label{thm:schur-common-cutoff-vector}
\lean{TensorPower.exists_commonLabelCutoff_vector}
    \leanok
\uses{thm:schur-common-cutoff-mass}
Under the hypotheses of Theorem~\ref{thm:schur-common-cutoff-mass},
suppose that $\psi$ is invariant under simultaneous copy permutations and
\begin{align}
 q&:=\|Q\psi\|^2-\sum_{j=1}^m\|(I-C_j)\psi\|^2>0.
\end{align}
There is a unit vector $\varphi$, invariant under simultaneous copy
permutations, with $Q\varphi=\varphi$ and $C_j\varphi=\varphi$ for every
$j$, and
\begin{align}
 |\langle\varphi,Q\psi\rangle|^2&=\|CQ\psi\|^2\ge q.
\end{align}
No normalization of $\psi$ is required.
\end{theorem}
\begin{proof}
    \leanok
Set $v=CQ\psi$. Theorem~\ref{thm:schur-common-cutoff-mass} gives
$v\ne0$, so take $\varphi=v/\|v\|$. The commuting projection
identities imply the range conditions. Each subsystem central projection
commutes with simultaneous copy permutations, which preserves the symmetry.
Since $C$ is an orthogonal projection,
$\langle v,Q\psi\rangle=\langle v,v\rangle$; normalization gives the overlap.
\end{proof}

\begin{theorem}[Preservation of the selected Bell range]
\label{thm:schur-common-cutoff-bell}
\lean{TensorPower.exists_commonLabelCutoff_bell_vector}
\leanok
\uses{thm:schur-common-cutoff-vector}
Let $\psi_0\in\mathbb C^A\otimes\mathbb C^B$, where $A$ and $B$ are
finite sets. Select spectral indices $E$ of its actual $A$-marginal,
and suppose their spectral mass $z_E$ is positive. Let
$J_E:\mathbb C^E\to\mathbb C^A$ be the corresponding eigenvector
isometry, and set
\begin{align}
 \beta_E&=|E|^{-1/2}\sum_{e\in E}J_E|e\rangle\otimes|e\rangle,
 &\Pi_E&=|\beta_E\rangle\langle\beta_E|.
\end{align}
Retain all the hypotheses on $F,k,\mathcal B,R,\lambda,\psi$ and $q$
of Theorem~\ref{thm:schur-common-cutoff-vector}. Then there is a unit
vector $\Phi$ in $(\mathbb C^A\otimes\mathbb C^E)^{\otimes k}$ tensored
with the exterior copy space, such that
\begin{align}
 (\Pi_E^{\otimes k}\otimes I)\Phi&=\Phi,
 &(I\otimes Q)\Phi&=\Phi,\\
 (I\otimes C_j)\Phi&=\Phi &&(1\le j\le m),\\
 \bigl|\langle\Phi,\beta_E^{\otimes k}\otimes Q\psi\rangle\bigr|^2
 &=\|CQ\psi\|^2.
\end{align}
After the canonical identification of copies of a product with the product
of copy spaces, $\Phi$ is invariant under simultaneous copy permutations.
The original vector $\psi_0$ need not be normalized.
\end{theorem}
\begin{proof}
\leanok
Take $\Phi=\beta_E^{\otimes k}\otimes\varphi$, where $\varphi$ is given
by Theorem~\ref{thm:schur-common-cutoff-vector}. Positive spectral mass
implies $E\ne\varnothing$, hence $\|\beta_E\|=1$. Tensor products give
the range conditions, the symmetry, and the exact overlap identity.
\end{proof}

\section{Physical regions and the prescribed compressed-state thresholds}
\label{sec:compressed_physical_regions}

Let $V$ be a finite set of physical sites. Fix $X\subseteq V$ and the augmented site
set $\widehat V=\{r\}\sqcup X^c$ of
Definition~\ref{def:compressed_site_family}. These identifications express
the nested exterior regions and their thresholds from
\cite[Section~7]{OpenAI2026AreaLaw} in the same site coordinates as the
actual compressed vector.

\begin{definition}[Extracting the physical sites]
\label{def:extract_compressed_physical_region}
\lean{FiniteProduct.extractExteriorRegion}
\leanok
For a region $D\subseteq\widehat V$, let
\begin{align}
 \operatorname{phys}_X(D)&=\{v\in X^c: v\text{ occurs as a physical site in }D\},
\end{align}
regarded as a subset of the original site set $V$. The auxiliary site
contributes no element to this set.
\end{definition}

\begin{theorem}[Recovery of an original exterior region]
\label{thm:extract_lifted_physical_region}
\lean{FiniteProduct.extractExteriorRegion_exteriorRegion}
\leanok
\uses{def:extract_compressed_physical_region,def:compressed_site_family}
If $B\subseteq V$ is disjoint from $X$, then
\begin{align}
 \operatorname{phys}_X(\widehat B)&=B.
\end{align}
\end{theorem}
\begin{proof}
\leanok
A physical site in $\widehat B$ belongs to $B$. Conversely, disjointness
places every site of $B$ in $X^c$, so that it occurs in $\widehat B$.
\end{proof}

\begin{theorem}[Inclusion, nesting, and auxiliary disjointness]
\label{thm:lifted_physical_regions_nested}
\lean{FiniteProduct.monotone_exteriorRegion,FiniteProduct.exteriorRegion_disjoint_auxiliary,FiniteProduct.pairwise_exteriorRegion}
\leanok
\uses{def:compressed_site_family}
For arbitrary physical regions $B,C\subseteq V$,
\begin{align}
 B\subseteq C&\ \Longrightarrow\ \widehat B\subseteq\widehat C,
 &\widehat B\cap\{r\}&=\varnothing.
\end{align}
Consequently, a finite list whose regions are pairwise comparable by
inclusion remains pairwise comparable after lifting. The list may contain
repeated regions.
\end{theorem}
\begin{proof}
\leanok
Intersecting with $X^c$ preserves inclusion, and the lift contains only
physical sites. Apply preservation of inclusion to each comparison in the
original list.
\end{proof}

\begin{definition}[The threshold determined by the original state]
\label{def:compressed_regional_threshold}
\lean{FiniteProduct.compressedRegionalThreshold}
\leanok
\uses{def:extract_compressed_physical_region}
Let the local physical coordinate sets be finite. Given an original vector
$\Omega$, a spectral selection $E$ on $X$ with actual mass $z$, and a
copy number $k\ge0$, define for $D\subseteq\widehat V$
\begin{align}
 a_k(D)&=k\left(\frac{S(\rho^\Omega_{\operatorname{phys}_X(D)})}{z}+1\right).
\end{align}
Here the formula is extended to $z=0$ by setting $a_k(D)=k$. Positive
selected mass is required when this threshold is used in the exponential
cutoff estimate.
\end{definition}

\begin{theorem}[The threshold on a lifted physical region]
\label{thm:compressed_regional_threshold_lifted}
\lean{FiniteProduct.compressedRegionalThreshold_exteriorRegion}
\leanok
\uses{def:compressed_regional_threshold,thm:extract_lifted_physical_region}
For every $B$ disjoint from $X$,
\begin{align}
 a_k(\widehat B)&=k\left(\frac{S(\rho_B^\Omega)}{z}+1\right),
\end{align}
with the same convention at $z=0$. No normalization of $\Omega$ is needed
for this identity.
\end{theorem}
\begin{proof}
\leanok
Substitute $\operatorname{phys}_X(\widehat B)=B$ in the definition.
\end{proof}

\section{Compressed physical and auxiliary sites}
\label{sec:compressed_site_coordinates}

\begin{definition}[The auxiliary site and exterior physical regions]
    \label{def:compressed_site_family}
    \lean{FiniteProduct.compressedSiteSpace,FiniteProduct.exteriorRegion}
    \leanok
    Let $V$ be a finite set of physical sites, with local coordinate
    sets $I_v$, and let $X\subseteq V$. For an auxiliary coordinate
    set $R$, define
    \begin{align}
        \widehat V&=\{r\}\sqcup X^c,&
        \widehat I_r&=R,&
        \widehat I_v&=I_v\quad(v\in X^c).
    \end{align}
    A physical region $B\subseteq V$ becomes
    $\widehat B=B\cap X^c$, included in $\widehat V$ and excluding the
    auxiliary site. These are the exterior coordinates in
    \cite[\texttt{07-comparators.tex}, lines 240--247 and
    283--308]{OpenAI2026AreaLaw}. They require no choice of a vector or
    spectral subspace, and no coordinate set is required to be nonempty.
\end{definition}

\section{Copy coordinates of the actual auxiliary factor}
\label{sec:compressed_typical_auxiliary_coordinates}

Use the augmented site family $\widehat V=\{r\}\sqcup X^c$ and local
coordinates of Definition~\ref{def:compressed_site_family}. The
auxiliary region is precisely $\{r\}$. The following identities use the
physical-complement copies first and the auxiliary copies second, as in
the Schmidt-selected vector of \cite[Section~7]{OpenAI2026AreaLaw}.

\begin{definition}[Collecting the exterior copies]
\label{def:compressed_exterior_copy_coordinates}
\lean{TensorPower.exteriorCopyEquiv}
\leanok
\uses{def:compressed_site_family}
For $k\ge0$, define the bijection
\begin{align}
 e_k:\left(\prod_{f\in\widehat V}\widehat I_f\right)^k
  &\longrightarrow\Omega_{X^c}^k\times A^k,\\
 e_k((x_j)_j)&=((x_j|_{X^c})_j,(x_j(r))_j).
\end{align}
Its inverse places each physical configuration and auxiliary index in the
same copy.
\end{definition}

\begin{theorem}[The auxiliary subsystem action]
\label{thm:compressed_auxiliary_copy_action}
\lean{TensorPower.exteriorCopyEquiv_subsystemPerm_none}
\leanok
\uses{def:compressed_exterior_copy_coordinates}
Let $\sigma\in\mathfrak S_k$ act only on the copies of the auxiliary
site. If $e_k(x)=(b,a)$, then
\begin{align}
 e_k(\sigma_{\{r\}}x)&=(b,(a_{\sigma^{-1}(j)})_j).
\end{align}
Thus the physical-complement coordinates are fixed, and the auxiliary
coordinates carry the ordinary copy action.
\end{theorem}
\begin{proof}
\leanok
Every physical site lies outside $\{r\}$, while the auxiliary site lies
inside it. Substitution in the definition of the subsystem action gives
the asserted action formula.
\end{proof}

\begin{theorem}[The same tensor-power vector in product coordinates]
\label{thm:compressed_exterior_product_coordinates}
\lean{TensorPower.exteriorCopyEquiv_prod}
\leanok
\uses{def:compressed_exterior_copy_coordinates}
For any function $\chi:A\times\Omega_{X^c}\to\mathbb C$, set
$\chi_{\mathrm{site}}=\chi\circ e$. Then for every
$(b,a)\in\Omega_{X^c}^k\times A^k$,
\begin{align}
 \prod_{j=0}^{k-1}\chi_{\mathrm{site}}\bigl((e_k^{-1}(b,a))_j\bigr)
 &=\prod_{j=0}^{k-1}\chi(a_j,b_j).
\end{align}
\end{theorem}
\begin{proof}
\leanok
In each copy, the inverse coordinate identification has auxiliary value
$a_j$ and physical-complement configuration $b_j$. The factors on the two
sides are therefore equal term by term.
\end{proof}

\begin{theorem}[Simultaneous-copy symmetry]
\label{thm:compressed_site_product_symmetry}
\lean{TensorPower.compressedSite_prod_mem_symmetricSubspace}
\leanok
\uses{def:compressed_site_family}
Assume that $A$ and the local physical coordinate sets are finite. For
every function $\chi:A\times\Omega_{X^c}\to\mathbb C$, the vector
\begin{align}
 \Psi_k((x_j)_j)&=\prod_{j=0}^{k-1}\chi_{\mathrm{site}}(x_j)
\end{align}
is fixed by every simultaneous permutation of its $k$ copies. No
normalization assumption is needed.
\end{theorem}
\begin{proof}
\leanok
The permutation changes only the order of the factors in the finite
product.
\end{proof}

For the actual compressed spectral vector, the product in
Theorem~\ref{thm:compressed_exterior_product_coordinates} uses exactly
the selected coordinate $a_j\in\{0,\ldots,|E|-1\}$ and the original
physical-complement configuration $b_j$. It therefore identifies the
given vector without making another choice of a spectral set or a
Schur-label sequence.

\section{Common cutoffs for the actual compressed state}
\label{sec:compressed_typical_common_cutoffs}

Let $V$ be a finite set of sites with finite coordinate sets, let
$\Omega$ be a unit physical vector, and fix $X\subseteq V$.
Select spectral indices $E$ of its actual marginal $\rho_X^\Omega$,
with $d=|E|$ and mass $z=\sum_{i\in E}p_i>0$. In the fixed enumeration
of $E$, let $J_E:\mathbb C^d\to\mathcal H_X$ be the eigenvector isometry.
Place the vector $\chi=z^{-1/2}(J_E^*\otimes I)\Omega$ at the augmented
sites $\widehat V=\{r\}\sqcup X^c$ to obtain $\chi_{\mathrm{site}}$,
and put $\Psi_k=\chi_{\mathrm{site}}^{\otimes k}$.
These are the actual selected coordinates of
\cite[Section~7, lines~240--247 and 332--354]{OpenAI2026AreaLaw}.

Fix a finite list $\mathcal B=(B_1,\ldots,B_m)$ of regions disjoint from
$X$, pairwise comparable by inclusion. The list may be empty or contain
repetitions. For any prescribed sequence of irreducible labels $\lambda_k$
of $\mathfrak S_k$, define
\begin{align}
 Q_k&=P_{\{r\}}^{\lambda_k},
 &C_{j,k}&=C_{\widehat B_j}\!\left(k\left(S(\rho_{B_j}^\Omega)/z+1\right)\right),\\
 C_k&=\prod_{j=1}^m C_{j,k},
 &q_k&=\|Q_k\Psi_k\|^2,\\
 \ell_k&=\sum_{j=1}^m\|(I-C_{j,k})\Psi_k\|^2,
 &r_k&=\|C_kQ_k\Psi_k\|^2.
 \label{eq:compressed_common_masses}
\end{align}
The cutoff thresholds use the original regional entropies and the same
selected mass $z$. Each product follows the order of the prescribed list.

\begin{theorem}[Retention of the same auxiliary-label mass]
\label{thm:compressed_typical_common_retained_mass}
\lean{FiniteProduct.eventually_commonLabelCutoff_compressedTypicalSite_retains_mass}
\leanok
\uses{thm:compressed_typical_regional_exponential_tail,thm:half_mass_after_finite_exponential_losses,thm:schur-common-cutoff-mass,thm:lifted_physical_regions_nested,thm:compressed_regional_threshold_lifted}
Assume that $q_k\ge 1/(2(k+1)^{d^2})$ for all sufficiently large $k$.
Then, for all sufficiently large $k$,
\begin{align}
 q_k-\ell_k&\ge q_k/2,
 &q_k-\ell_k&\ge\frac{1}{4(k+1)^{d^2}}>0,\\
 r_k&\ge q_k/2,
 &r_k&\ge\frac{1}{4(k+1)^{d^2}}.
\end{align}
The label sequence is the prescribed sequence in
(\ref{eq:compressed_common_masses}).
\end{theorem}
\begin{proof}\leanok
For each fixed list position, Theorem~\ref{thm:compressed_typical_regional_exponential_tail}
bounds its loss by $e^{-c_jk}$ for some $c_j>0$ and every $k\ge1$.
Theorem~\ref{thm:half_mass_after_finite_exponential_losses} gives the two
bounds on $q_k-\ell_k$, and its lower bound is strictly positive.
The lifted regions remain nested and avoid the auxiliary site by
Theorem~\ref{thm:lifted_physical_regions_nested}. Thus the finite projection
union bound gives $q_k-\ell_k\le r_k$, proving the remaining inequalities.
\end{proof}

\begin{theorem}[A Bell vector with the same common-cutoff mass]
\label{thm:compressed_typical_common_bell_vector}
\lean{FiniteProduct.eventually_exists_commonLabelCutoff_compressedTypical_bell_vector}
\leanok
\uses{thm:compressed_typical_common_retained_mass,thm:compressed_site_product_symmetry,thm:schur-common-cutoff-bell}
Assume the same eventual lower bound
$q_k\ge 1/(2(k+1)^{d^2})$. Let $\beta_E$ be the actual selected Bell vector
on $\mathcal H_X\otimes\mathbb C^E$, with projection
$\Pi_E=|\beta_E\rangle\langle\beta_E|$, and put
$\mathrm{post}_k=\beta_E^{\otimes k}\otimes\Psi_k$.
For all sufficiently large $k$, there is a unit vector $\varphi_k$ such that
\begin{align}
 (\Pi_E^{\otimes k}\otimes I)\varphi_k&=\varphi_k,
 &(I\otimes Q_k)\varphi_k&=\varphi_k,\\
 (I\otimes C_{j,k})\varphi_k&=\varphi_k &&(1\le j\le m),\\
 \bigl|\langle\varphi_k,(I\otimes Q_k)\mathrm{post}_k\rangle\bigr|^2
 &=r_k,\\
 r_k&\ge q_k/2,
 &r_k&\ge\frac{1}{4(k+1)^{d^2}}.
\end{align}
The vector $\varphi_k$ is invariant under simultaneous copy permutations,
using the canonical identification of copies of a product with the product
of copy spaces. The left Bell coordinate is indexed by $E$ and the right
auxiliary coordinate by its fixed enumeration $\{0,\ldots,d-1\}$.
The eventual threshold may depend on the original state, selection,
region list, and prescribed label sequence.
\end{theorem}
\begin{proof}\leanok
Theorem~\ref{thm:compressed_typical_common_retained_mass} makes
$q_k-\ell_k$ eventually positive and gives the bounds on $r_k$.
The vector $\Psi_k$ is symmetric by
Theorem~\ref{thm:compressed_site_product_symmetry}. Apply
Theorem~\ref{thm:schur-common-cutoff-bell} to this actual exterior vector,
with the original selected Bell vector and the same label $\lambda_k$.
Its exact overlap identity gives $r_k$ from
(\ref{eq:compressed_common_masses}).
\end{proof}

\section{The same auxiliary-label mass after regrouping}
\label{sec:compressed_typical_label_mass}

The following coordinate identity identifies the auxiliary projection in
the site description with the projection used to choose the label in
\cite[Section~7]{OpenAI2026AreaLaw}.



\begin{theorem}[The actual auxiliary isotypic projection]
\label{thm:compressed_site_label_projection_coordinates}
\lean{TensorPower.labelProj_compressedSite_submatrix}
\leanok
\uses{thm:group_algebra_product_coordinates,thm:compressed_auxiliary_copy_action}
Use the augmented site family $\widehat V=\{r\}\sqcup X^c$, with
finite auxiliary coordinate set $A$ and finite local physical coordinate
sets. Let $T_k$ be the unitary induced by the bijection $e_k$ of
Definition~\ref{def:compressed_exterior_copy_coordinates}. For every
$k\ge0$ and every irreducible label $\lambda$ of $\mathfrak S_k$,
\begin{align}
 T_k\Pi_{\{r\},\lambda}T_k^*
 &=I_{\Omega_{X^c}^k}\otimes\Pi^{A,k}_\lambda,
\end{align}
where the left projection is for the actual auxiliary subsystem action
and the right projection is for the ordinary copy action on $A^k$.
The label $\lambda$ is unchanged.
\end{theorem}
\begin{proof}
\leanok
The auxiliary copy-action identity supplies the product equivariance.
Apply Theorem~\ref{thm:group_algebra_product_coordinates} to the central idempotent of the given
irreducible label.
\end{proof}

\begin{theorem}[Equality of the auxiliary projected masses]
\label{thm:compressed_site_label_mass_coordinates}
\lean{TensorPower.norm_sq_labelProj_compressedSite_prod,TensorPower.norm_sq_labelProj_compressedTypicalSite_prod}
\leanok
\uses{thm:compressed_site_label_projection_coordinates,thm:compressed_exterior_product_coordinates}
For any function $\chi:A\times\Omega_{X^c}\to\mathbb C$, define
\begin{align}
 \Psi_k((x_j)_j)&=\prod_{j=0}^{k-1}\chi(e(x_j)),
 &w_k(b,a)&=\prod_{j=0}^{k-1}\chi(a_j,b_j).
\end{align}
For every irreducible label $\lambda$ of $\mathfrak S_k$,
\begin{align}
 \|\Pi_{\{r\},\lambda}\Psi_k\|^2
 &=\|(I_{\Omega_{X^c}^k}\otimes\Pi^{A,k}_\lambda)w_k\|^2.
\end{align}
In particular, take $A=\{0,\ldots,|E|-1\}$ and take $\chi$ to be
the actual compressed spectral vector of an original physical vector
$\Omega$, with the selected indices in their fixed order. The right
side is then exactly the auxiliary-label mass of that same compressed
vector in physical-complement-first coordinates. The identity applies
to every member of an already chosen label sequence. It includes zero
copies and zero selected mass and requires no normalization or symmetry
assumption.
\end{theorem}
\begin{proof}
\leanok
The product-coordinate identity gives $T_k\Psi_k=w_k$. Combine this
with Theorem~\ref{thm:compressed_site_label_projection_coordinates} and use preservation of norm by
$T_k$. Substitution of the actual compressed spectral vector gives the
last assertion without making another selection of a vector or a label.
\end{proof}

\section{Regional marginals after spectral compression}
\label{sec:compressed_typical_regional_marginal}

Let $F$ be a finite set of sites with finite local coordinate sets, and
let $\Omega$ be any vector in the resulting tensor product. Fix disjoint
regions $X,B\subseteq F$. View $\Omega$ in
$\mathcal H_X\otimes\mathcal H_{X^c}$, and choose a set $E$ of
eigenvalue indices of its marginal on $X$. Let
$J_E:\mathbb C^{|E|}\to\mathcal H_X$ have the corresponding
orthonormal eigenvectors as columns, ordered by the fixed enumeration of
$E$, and set
\begin{align}
 P_E&=J_EJ_E^*,
 &z&=\operatorname{Tr}(P_E\rho_X),\\
 \psi&=z^{-1/2}(P_E\otimes I)\Omega,
 &\chi&=z^{-1/2}(J_E^*\otimes I)\Omega.
\end{align}
The vector $\psi$ is regarded in the original physical coordinates,
whereas $\chi\in\mathbb C^{|E|}\otimes\mathcal H_{X^c}$.
When $z=0$, set $\psi=\chi=0$. This is the
selected spectral compression used in
\cite[Section~7]{OpenAI2026AreaLaw}.

\begin{theorem}[The marginal on a disjoint region is preserved]
\label{thm:compressed_typical_regional_marginal}
\lean{FiniteProduct.partialTraceLeft_compressedTypicalPureState_region}
\leanok
\uses{thm:compressed_typical_vector}
Use the canonical physical coordinate split
$\mathcal H_{X^c}=\mathcal H_B\otimes\mathcal H_{(X\cup B)^c}$
to write $\chi(r,b,c)$. Place the discarded coordinates first by setting
$\tau((r,c),b)=\chi(r,b,c)$. Then
\begin{align}
 \operatorname{Tr}_{\mathbb C^{|E|}\otimes\mathcal H_{(X\cup B)^c}}
       |\tau\rangle\langle\tau|
 &=\operatorname{Tr}_{\mathcal H_{B^c}}|\psi\rangle\langle\psi|.
\end{align}
There is no normalization assumption on $\Omega$ and no positive-mass
assumption on $E$.
\end{theorem}
\begin{proof}
\leanok
Embedding the selected coordinate space by $J_E$ takes $\chi$ to
$\psi$. Since $J_E$ is an isometry, this embedding preserves the entire
marginal on $X^c$. Trace that marginal further over $(X\cup B)^c$.
In coordinates, this amounts to replacing the sum over the selected
index $r$ by the sum over the original $X$ coordinates and then
interchanging the two finite sums. The resulting matrix is precisely the
regional marginal of the same selected vector $\psi$.
\end{proof}

\section{The actual compressed state's regional cutoff losses}
\label{sec:compressed_typical_regional_tail}

Let $V$ be a finite set of sites with finite local coordinate sets, let
$\Omega$ be a physical vector, and fix a region $X\subseteq V$.
Select eigenvalue indices $E$ of the actual marginal $\rho_X^\Omega$,
with selected mass $z=\operatorname{Tr}(P_E\rho_X^\Omega)>0$.
Let $J_E$ have those orthonormal eigenvectors as columns in their fixed
order. In the $X,X^c$ coordinates, write
\begin{align}
 \chi&=z^{-1/2}(J_E^*\otimes I)\Omega.
\end{align}
Place its selected coordinate at the auxiliary site $r$ and its remaining
coordinates at their original sites in $X^c$, obtaining the actual vector
$\chi_{\mathrm{site}}$ on $\widehat V=\{r\}\sqcup X^c$.
These are the selected state and mass of
\cite[Section~7]{OpenAI2026AreaLaw}.

\begin{theorem}[Normalization in the actual site coordinates]
\label{thm:compressed_typical_site_norm}
\lean{FiniteProduct.norm_compressedTypicalSiteState}
\leanok
\uses{def:compressed_site_family,thm:compressed_typical_vector}
The actual compressed vector satisfies
\begin{align}
 \|\chi_{\mathrm{site}}\|&=1.
\end{align}
Only $z>0$ is required; the original vector $\Omega$ need not be unit.
\end{theorem}
\begin{proof}
\leanok
The coordinate identification first separates the auxiliary site and then
identifies its ordered indices with $E$. It is an isometry from the
selected coordinate space to the site description. The selected compressed
vector has squared norm $z^{-1}z=1$.
\end{proof}

\begin{theorem}[The compressed regional entropy bound]
\label{thm:compressed_typical_site_entropy}
\lean{FiniteProduct.entropy_compressedTypicalSiteState_le}
\leanok
\uses{thm:compressed_site_actual_regional_marginal}
Assume in addition that $\|\Omega\|=1$. For every physical region
$B$ disjoint from $X$, regard $B$ as the exterior-site region
$\widehat B$. Then
\begin{align}
 S\!\left(\rho_{\widehat B}^{\chi_{\mathrm{site}}}\right)
 &\le\frac{S(\rho_B^\Omega)}{z}.
\end{align}
\end{theorem}
\begin{proof}
\leanok
The actual compressed regional marginal, after the canonical coordinate
identification, equals the marginal of the actual selected physical
vector. Entropy is unchanged by this identification. The selected-state
entropy inequality bounds the latter entropy by $S(\rho_B^\Omega)/z$.
The mass $z$ is the same original spectral-selection mass throughout.
\end{proof}

\begin{theorem}[Exponential loss of an actual regional Schur cutoff]
\label{thm:compressed_typical_regional_exponential_tail}
\lean{FiniteProduct.exists_norm_sq_one_sub_labelCutoff_compressedTypicalSite_prod_le_exp}
\leanok
\uses{thm:compressed_typical_site_norm,thm:compressed_typical_site_entropy,thm:iid_schur_label_exponential,thm:subsystem_label_cutoff_regional_trace}
Assume $\|\Omega\|=1$ and $X\cap B=\varnothing$. Let
$P_{\widehat B}^{(k)}(a)$ be the closed central-label cutoff
$\log\dim\lambda\le a$ for the actual subsystem copy action on
$\widehat B$. There is a constant $c>0$ such that for every integer
$k\ge1$,
\begin{align}
 \left\|\left(I-P_{\widehat B}^{(k)}
      \!\left(k\left(\frac{S(\rho_B^\Omega)}{z}+1\right)\right)\right)
      \chi_{\mathrm{site}}^{\otimes k}\right\|^2
 &\le e^{-ck}.
\end{align}
The rate may depend on the fixed original state, selection, and region.
The region may be empty and its marginal may be singular.
\end{theorem}
\begin{proof}
\leanok
The actual regional marginal of $\chi_{\mathrm{site}}$ is positive
and has trace one by Theorem~\ref{thm:compressed_typical_site_norm}.
Its entropy is bounded by Theorem~\ref{thm:compressed_typical_site_entropy}.
Apply Theorem~\ref{thm:iid_schur_label_exponential} with
$s=S(\rho_B^\Omega)/z$, then use
Theorem~\ref{thm:subsystem_label_cutoff_regional_trace} to identify the trace
tail with the squared cutoff failure norm of $\chi_{\mathrm{site}}^{\otimes k}$.
\end{proof}

\section{The compressed vector as a family of physical sites}
\label{sec:compressed_typical_site_coordinates}

Let $V$ be a finite set of physical sites, let $I_v$ be its local
coordinate sets, and fix $X\subseteq V$. Write
$\Omega_D=\prod_{v\in D}I_v$ for the configurations of a region $D$.
The selected spectral coordinates can be regarded as one additional site,
as in the rough comparison of \cite[Section~7]{OpenAI2026AreaLaw}.

Use the augmented family $\widehat V=\{r\}\sqcup X^c$ of
Definition~\ref{def:compressed_site_family}, with auxiliary coordinate set
$A$, and write $\widehat B=B\cap X^c$ for its physical regions.
The canonical coordinate bijection is
\begin{align}
 e:\prod_{f\in\widehat V}\widehat I_f
   &\longrightarrow A\times\Omega_{X^c},
 &e(x)&=(x_r,(x_v)_{v\in X^c}).
\end{align}

\begin{definition}[Regional and discarded coordinates]
\label{def:compressed_site_region_coordinates}
\lean{FiniteProduct.exteriorRegionEquiv,FiniteProduct.exteriorComplementEquiv}
\leanok
\uses{def:compressed_site_family}
For every $B\subseteq V$, evaluation of the auxiliary and physical
coordinates gives a bijection
\begin{align}
 e_C:\prod_{f\in\widehat B^c}\widehat I_f
   &\longrightarrow A\times\Omega_{(X\cup B)^c}.
\end{align}
If $X\cap B=\varnothing$, evaluation of the physical coordinates also
gives a bijection
\begin{align}
 e_B:\prod_{f\in\widehat B}\widehat I_f
   &\longrightarrow\Omega_B.
\end{align}
The inverses assign these values to the same auxiliary and physical sites.
No local coordinate set is required to be nonempty or finite for these
bijections.
\end{definition}

\begin{theorem}[Compatibility of the regional split]
\label{thm:compressed_site_region_split}
\lean{FiniteProduct.piOptionEquivProd_split_exteriorRegion}
\leanok
\uses{def:compressed_site_region_coordinates}
Suppose $X\cap B=\varnothing$. Let $b$ and $c$ be configurations on
$\widehat B$ and $\widehat B^c$, respectively, and write $e_C(c)=(a,d)$.
Joining $b$ and $c$ into one exterior-site configuration and then applying
$e$ gives
\begin{align}
 e(b\cup c)&=(a,e_B(b)\cup d).
\end{align}
On the right, the union joins the physical configurations on $B$ and
$(X\cup B)^c$ to a configuration on $X^c$.
\end{theorem}
\begin{proof}
\leanok
At the auxiliary site, both sides have value $a$. At a physical site in
$B$, both have the value prescribed by $b$; at every other site of $X^c$,
both have the value prescribed by $d$.
\end{proof}

\begin{theorem}[The actual selected marginal in exterior-site coordinates]
\label{thm:compressed_site_actual_regional_marginal}
\lean{FiniteProduct.reducedPure_compressedTypicalSiteState}
\leanok
\uses{thm:compressed_site_region_split,thm:compressed_typical_regional_marginal}
Assume now that the local coordinate sets are finite. Let $\Omega$ be
any original physical vector, and let $E$ select eigenvalue indices of its
marginal on $X$. Let $\psi$ be its actual selected vector and let
$\chi\in\mathbb C^{|E|}\otimes\mathbb C^{\Omega_{X^c}}$ be its
actual compressed vector, with the selected indices in their fixed order.
Take $A=\{0,\ldots,|E|-1\}$ and define
$\chi_{\mathrm{site}}(x)=\chi(e(x))$.

If $X\cap B=\varnothing$ and $U_B$ is the unitary coordinate
identification induced by $e_B$, then
\begin{align}
 U_B\operatorname{Tr}_{\widehat B^c}
   |\chi_{\mathrm{site}}\rangle\langle\chi_{\mathrm{site}}|U_B^*
 &=\operatorname{Tr}_{B^c}|\psi\rangle\langle\psi|.
\end{align}
The identity does not require normalization of $\Omega$ or positive
selected mass.
\end{theorem}
\begin{proof}
\leanok
The split identity identifies the discarded coordinates with the selected
index and the physical complement of $X\cup B$. Reindex the finite
partial-trace sum by $e_C$, and apply preservation of the regional marginal
under the actual spectral compression.
\end{proof}

\section{Permutation representations in product coordinates}
\label{sec:group_algebra_product_coordinates}

The coordinate identity below applies to the regional copy actions in
\cite[Section~5, lines~14--27 and 73--80]{OpenAI2026AreaLaw}.

\begin{theorem}[Group-algebra operators in product coordinates]
\label{thm:group_algebra_product_coordinates}
\lean{PermutationRepresentation.groupAlgebraRep_submatrix_of_prod_action}
    \leanok
Let a group $G$ act by permutations on finite sets $Z$ and $B$, with
actions $\alpha$ and $\beta$. Let $A$ be any set and
$e:Z\to A\times B$ a bijection satisfying
\begin{align}
 e(\alpha(g)z)&=((ez)_A,\beta(g)(ez)_B)
\end{align}
for every $g\in G$ and $z\in Z$. Let $\pi_\alpha$ and $\pi_\beta$ be
the linear extensions of the two permutation representations to
$\mathbb C[G]$. For every $h\in\mathbb C[G]$, $a,a'\in A$, and
$b,b'\in B$,
\begin{align}
 \bigl(\pi_\alpha(h)\bigr)_{e^{-1}(a,b),e^{-1}(a',b')}
   &=\delta_{a,a'}\bigl(\pi_\beta(h)\bigr)_{b,b'}.
\end{align}
Equivalently, the matrix in product coordinates is
$I_A\otimes\pi_\beta(h)$. Neither $G$ nor $A$ is required to be finite.
\end{theorem}
\begin{proof}
    \leanok
For a single group element, the permutation-matrix entry at $(x,y)$ is
one precisely when $\alpha(g)y=x$. Equivariance expresses this as
equality of the first coordinates and the corresponding permutation
condition on the second coordinates. This gives the identity for group
elements. Extend it by linearity over the finite support of $h$.
\end{proof}

\section{Exponential upper tails of the actual copy labels}
\label{sec:iid_schur_label_exponential}

For a finite density matrix $\rho$, let $\Pi^{(k)}_\lambda$ be the
central isotypic projections of the permutation action on its $k$ copies,
and define the closed label cutoff
\begin{align}
 P_k(a)&=\sum_{\lambda:\,\log\dim\lambda\le a}\Pi^{(k)}_\lambda.
\end{align}
The following consequence of the supported comparison with surprisal is
the fixed-system exponential estimate used in
\cite[Section~7]{OpenAI2026AreaLaw}.

\begin{theorem}[Independent-copy Schur tails above an entropy bound]
\label{thm:iid_schur_label_exponential}
\lean{Matrix.PosSemidef.exists_re_trace_one_sub_labelCutoff_finKronecker_le_exp}
\leanok
\uses{thm:schur_label_tail_surprisal,thm:iid_surprisal_exponential_density,thm:surprisal_functional_calculus}
Let $\rho\ge0$ be a finite matrix with $\operatorname{Tr}\rho=1$,
and let $s\in\mathbb R$ satisfy $S(\rho)\le s$. There is a constant
$c>0$ such that, for every integer $k\ge1$,
\begin{align}
 \operatorname{Re}\operatorname{Tr}\!\left(
   \rho^{\otimes k}\bigl(I-P_k(k(s+1))\bigr)\right)
 &\le e^{-ck}.
\end{align}
Singular density matrices are included. The permutation invariance is that of the actual tensor power.
\end{theorem}
\begin{proof}
\leanok
The actual tensor power $\rho_k=\rho^{\otimes k}$ is positive, has
trace one, and commutes with all copy permutations. The supported Schur
comparison therefore bounds its removed label mass at threshold $a$ by
\begin{align}
 \operatorname{Re}\operatorname{Tr}\!\left(
   \rho_k\mathbf1_{(a,\infty)}(-\log\rho_k)\right),
\end{align}
where the logarithm is zero on the kernel. The functional-calculus
composition identity gives precisely this cutoff expression. Since
$k(S(\rho)+1)\le k(s+1)$, the strict upper tail at the latter threshold
is contained in that at the former. The eigenvalue weights are
nonnegative, so this inclusion also orders their trace masses. Apply the
exponential surprisal bound at $k(S(\rho)+1)$.
\end{proof}

\section{Spectral cutoffs of the label observable}
\label{sec:schur_label_cutoffs}

Let a finite group $G$ act by permutations on a finite set $X$, and write
$\pi^\lambda$ for the orthogonal projections onto the irreducible labels.
Set $d_\lambda=\dim\lambda$ and
$F=\sum_\lambda(\log d_\lambda)\pi^\lambda$.
The following statements express the supported tail comparison used in
\cite[Section~7, upper comparisons]{OpenAI2026AreaLaw}.

\begin{definition}[The closed label cutoff]
    \label{def:schur_label_cutoff}
    \lean{PermutationRepresentation.labelCutoff}
    \leanok
    \uses{def:label_projectors}
    For $a\in\mathbb R$, define
    \begin{align*}
        C(a) & =\sum_{\log d_\lambda\le a}\pi^\lambda.
    \end{align*}
\end{definition}

\begin{theorem}[The label cutoff is a spectral projection]
    \label{thm:schur_label_cutoff_spectral}
    \lean{PermutationRepresentation.labelCutoff_eq_spectralProjectionGE}
    \leanok
    \uses{def:schur_label_cutoff,lem:label_projectors,lem:orthogonal_resolution_cfc}
    For every $a\in\mathbb R$,
    \begin{align*}
        C(a) & =\mathbf1_{(-\infty,a]}(F)
        =\mathbf1_{[-a,\infty)}(-F).
    \end{align*}
    The cutoff includes labels exactly at the threshold.
\end{theorem}
\begin{proof}\leanok
    Compute the functional calculus in the orthogonal label decomposition.
    The coefficient on the $\lambda$ subspace is one precisely when
    $\log d_\lambda\le a$.
\end{proof}

\begin{theorem}[Label tails are bounded by surprisal tails]
    \label{thm:schur_label_tail_surprisal}
    \lean{PermutationRepresentation.re_trace_mul_one_sub_labelCutoff_le}
    \leanok
    \uses{def:schur_label_cutoff,thm:schur_label_surprisal}
    Let $\rho$ be a positive semidefinite matrix on $\mathbb C^X$ with
    $\operatorname{Tr}\rho=1$, commuting with every permutation operator
    of the action. For every $a\in\mathbb R$,
    \begin{align*}
        \operatorname{Re}\operatorname{Tr}\bigl(\rho(1-C(a))\bigr)
         & \le \operatorname{Re}\operatorname{Tr}\bigl(
        \rho\mathbf1_{(a,\infty)}(-\log\rho)\bigr).
    \end{align*}
    The logarithm is assigned value zero on the kernel of $\rho$.
\end{theorem}
\begin{proof}\leanok
    Use the joint orthogonal decomposition for $\rho$ and the label
    projectors. On a nonzero joint subspace with eigenvalue $u>0$ and label
    $\lambda$, the trace-one condition gives $u d_\lambda\le1$.
    Consequently $\log d_\lambda\le-\log u$, so the event
    $\log d_\lambda>a$ is contained in $-\log u>a$.
    Multiply this scalar comparison by $u$ and sum the traces of the joint
    projections. Subspaces with $u=0$ contribute zero to both sides.
\end{proof}

\section{Regional copy coordinates and cutoff mass}
\label{sec:subsystem_label_coordinates}

Let $F$ be a finite set of sites, with a set of local coordinates $I_f$ at
each site. For $B\subseteq F$, write
$\Omega_B=\prod_{f\in B}I_f$ and $\Omega=\prod_{f\in F}I_f$.
These coordinates identify the subsystem permutation action used in the
rough comparison of \cite[Section~7]{OpenAI2026AreaLaw}.

\begin{definition}[Complementary and regional copy coordinates]
\label{def:subsystem_split_coordinates}
\lean{TensorPower.subsystemSplitEquiv}
    \leanok
For every $k\ge0$, define the bijection
\begin{align}
 e_B:\Omega^k&\longrightarrow\Omega_{B^c}^k\times\Omega_B^k,
 & e_B((x_j)_j)&=((x_j|_{B^c})_j,(x_j|_B)_j).
\end{align}
Its inverse joins the two restricted configurations in each copy.
\end{definition}

\begin{theorem}[The subsystem action in split coordinates]
\label{thm:subsystem_split_permutation}
\lean{TensorPower.subsystemSplitEquiv_subsystemPerm}
    \leanok
\uses{def:subsystem_split_coordinates}
Let $\sigma\in\mathfrak S_k$ act on the copies at sites in $B$, fixing
the coordinates at sites in $B^c$. If $e_B(x)=(c,r)$, then
\begin{align}
 e_B(\sigma_B x)&=(c,(r_{\sigma^{-1}(j)})_j).
\end{align}
\end{theorem}
\begin{proof}
    \leanok
Restriction to $B^c$ sees only fixed coordinates. Restriction to $B$ sees
exactly the copy permutation, with the inverse permutation in the
coordinate formula for the left action.
\end{proof}

Assume henceforth that every $I_f$ is finite. Let $T_B$ be the unitary
coordinate identification induced by $e_B$, so that
$(T_Bv)(c,r)=v(e_B^{-1}(c,r))$. Write $\Pi_{B,\lambda}$ for the central
isotypic projection of the subsystem action, and write
$\Pi^{\mathrm{reg}}_\lambda$ for that of the ordinary copy action on
$\mathbb C^{\Omega_B^k}$. Define the closed cutoffs
\begin{align}
 P_B(a)&=\sum_{\lambda:\,\log\dim\lambda\le a}\Pi_{B,\lambda},
 &P^{\mathrm{reg}}_B(a)
 &=\sum_{\lambda:\,\log\dim\lambda\le a}\Pi^{\mathrm{reg}}_\lambda.
\end{align}

\begin{theorem}[Regional Schur projections and cutoffs]
\label{thm:subsystem_label_projection_coordinates}
\lean{TensorPower.labelProj_subsystemPerm_submatrix,TensorPower.labelCutoff_subsystemPerm_submatrix}
    \leanok
\uses{thm:subsystem_split_permutation,thm:group_algebra_product_coordinates}
For every irreducible label $\lambda$ of $\mathfrak S_k$ and every
$a\in\mathbb R$,
\begin{align}
 T_B\Pi_{B,\lambda}T_B^*
   &=I_{\Omega_{B^c}^k}\otimes\Pi^{\mathrm{reg}}_\lambda,\\
 T_BP_B(a)T_B^*
   &=I_{\Omega_{B^c}^k}\otimes P^{\mathrm{reg}}_B(a).
\end{align}
\end{theorem}
\begin{proof}
    \leanok
Apply the group-algebra product-coordinate identity to the actual
subsystem action and the central idempotent of $\lambda$. Theorem~\ref{thm:subsystem_split_permutation} supplies its equivariance hypothesis. Summing the
resulting projection identities with $\log\dim\lambda\le a$ gives the
cutoff identity.
\end{proof}

\begin{theorem}[Cutoff failure as a regional trace]
\label{thm:subsystem_label_cutoff_regional_trace}
\lean{TensorPower.norm_sq_one_sub_labelCutoff_subsystemPerm_prod}
    \leanok
\uses{thm:subsystem_label_projection_coordinates}
For any vector $\psi\in\mathbb C^\Omega$, let
$\rho_B=\operatorname{Tr}_{B^c}|\psi\rangle\langle\psi|$ be its
actual regional marginal. For every $k\ge0$ and $a\in\mathbb R$,
\begin{align}
 \bigl\|(I-P_B(a))\psi^{\otimes k}\bigr\|^2
 &=\operatorname{Re}\operatorname{Tr}\!\left(
    \rho_B^{\otimes k}(I-P^{\mathrm{reg}}_B(a))\right).
\end{align}
No normalization of $\psi$ is required. The identity includes empty
regions, zero copies, and singular marginals.
\end{theorem}
\begin{proof}
    \leanok
In the split coordinates, the cutoff complement is the identity on the
complementary copies tensored with the regional cutoff complement. Both
are orthogonal projections. Hence its squared norm on the vector is its
expectation in the corresponding pure-state matrix. Tracing the
complementary copies gives $\rho_B^{\otimes k}$, which proves the
identity.
\end{proof}
